\chapter{刚体}
\section{刚体模型及其运动}
\subsection{刚体}
特点：刚体是一种理想模型。它可以看成由无数个质元组成的一种特殊的体系，在一定的外力的作用下，系统内任意两质元间的距离始终保持不变。

常见应用对象：外力作用下形变并不显著的物体。

刚体与质点系的关系：刚体是一种常用的质点系模型。质点系中的运动规律在刚体模型中同样适用，刚体还有一些特殊的运动规律。

\subsection{平动与转动}
刚体的一般运动可以分解为平动和转动。

\subsubsection{平动}
定义：刚体内任何一条给定的直线，在运动中它的方向始终保持不变。

运动特点：\begin{enumerate}
	\item 运动轨迹不一定是直线。
	\item 在任何时刻，各个质元的速度和加速度都是相等的。在任意一段时间内，所有质元的位移都是相同的。
	\item 刚体的平动运动可以用一个点的运动描述，通常用刚体的质心。
\end{enumerate}

\subsubsection{转动}
定义：刚体的各个质元在运动中都绕同一直线作圆周运动，这样的运动就叫做转动。如果转轴是固定不动的，就叫做定轴转动。

分类：定轴转动、定点转动……

\subsection{自由度}
定义：确定一个物体在空间的位置所需要的独立坐标的数目。

典型举例：\begin{itemize}
	\item 自由质点：3个自由度
	\item 平面运动质点：2个自由度
	\item 曲线运动质点：1个自由度
	\item 自由刚体：6个自由度（3个平动自由度，3个转动自由度）
	\item 刚体的定轴转动：1个自由度
	\item 刚体的定点转动：3个自由度
	\item 平面平行运动：3个自由度
\end{itemize}